% 2011 年天津卷理科第 10 题：三视图（正视图、侧视图、俯视图），单位 m。
% 来源 PDF：普通高考/2011/2011天津理.pdf，SHA-256 4d77f7437f1f7191d1fd3cd766c7c3014859dec8fdd7a4d8f4bca348c0ef06e5。
% 原页：第 1 页（A3 三栏）第 2 栏中部，书页码 705。
% 重绘依据：照原页——正视图为下底宽 3、高 1 的矩形上叠底宽 2、高 3 的三角形；
%   侧视图为下底宽 2、高 1 的矩形上叠底宽 2、高 3 的三角形；
%   俯视图为 3×2 的矩形加内切圆（半径 1）。
\begin{tikzpicture}[scale=0.52, >=stealth, semithick, font=\small]

  % ===== 正视图 =====
  \draw (0,0) rectangle (3,1);
  \draw (0.5,1) -- (1.5,4) -- (2.5,1);
  % 高 3
  \draw[<->] (3.7,1) -- (3.7,4);
  \node[right] at (3.7,2.5) {$3$};
  \draw (3,1) -- (3.85,1);
  \draw (3,4) -- (3.85,4);
  % 高 1
  \draw[<->] (3.7,0) -- (3.7,1);
  \node[right] at (3.7,0.5) {$1$};
  \draw (3,0) -- (3.85,0);
  % 宽 3
  \draw[<->] (0,-0.8) -- (3,-0.8);
  \node[below] at (1.5,-0.8) {$3$};
  \node[below] at (1.5,-1.5) {正视图};

  % ===== 侧视图 =====
  \begin{scope}[xshift=7.4cm]
    \draw (0,0) rectangle (2,1);
    \draw (0,1) -- (1,4) -- (2,1);
    \draw[<->] (2.7,1) -- (2.7,4);
    \node[right] at (2.7,2.5) {$3$};
    \draw (2,1) -- (2.85,1);
    \draw (2,4) -- (2.85,4);
    \draw[<->] (2.7,0) -- (2.7,1);
    \node[right] at (2.7,0.5) {$1$};
    \draw (2,0) -- (2.85,0);
    \draw[<->] (0,-0.8) -- (2,-0.8);
    \node[below] at (1,-0.8) {$2$};
    \node[below] at (1,-1.5) {侧视图};
  \end{scope}

  % ===== 俯视图 =====
  \begin{scope}[yshift=-6.6cm]
    \draw (0,0) rectangle (3,2);
    \draw (1.5,1) circle (1);
    \fill (1.5,1) circle (1.2pt);
    \draw[->] (1.5,1) -- (2.207,1.707);
    \node[above left] at (1.85,1.35) {$1$};
    % 高 2
    \draw[<->] (3.7,0) -- (3.7,2);
    \node[right] at (3.7,1) {$2$};
    \draw (3,0) -- (3.85,0);
    \draw (3,2) -- (3.85,2);
    % 宽 3
    \draw[<->] (0,-0.8) -- (3,-0.8);
    \node[below] at (1.5,-0.8) {$3$};
    \node[below] at (1.5,-1.5) {俯视图};
  \end{scope}

\end{tikzpicture}
